% =====================================================================
%  UNIVERSITÉ PARIS CITÉ — MASTER 2 MLDS
%  Human Activity Recognition using Unsupervised Learning
%  Author: Hamady GACKOU
%  Step 5 — Discussion & Report Integration
% =====================================================================

\documentclass[12pt,a4paper]{report}

% -------------------------------------------------------------
% PACKAGES
% -------------------------------------------------------------
\usepackage[utf8]{inputenc}
\usepackage[T1]{fontenc}
\usepackage[french,english]{babel}
\usepackage{geometry}
\usepackage{graphicx}
\usepackage{booktabs}
\usepackage{caption}
\usepackage{subcaption}
\usepackage{amsmath,amsfonts,amssymb}
\usepackage{hyperref}
\usepackage{setspace}
\usepackage{enumitem}
\usepackage{xcolor}
\usepackage{float}
\usepackage{titlesec}
\usepackage{lmodern}

% -------------------------------------------------------------
% PAGE STYLE
% -------------------------------------------------------------
\geometry{top=2.5cm,bottom=2.5cm,left=2.5cm,right=2.5cm}
\setstretch{1.25}
\hypersetup{
    colorlinks=true,
    linkcolor=blue,
    urlcolor=cyan,
    citecolor=teal
}
\titleformat{\chapter}[block]{\normalfont\huge\bfseries}{\thechapter.}{1em}{}
\titleformat{\section}[block]{\normalfont\Large\bfseries}{\thesection.}{0.75em}{}
\renewcommand{\figurename}{Figure}
\renewcommand{\tablename}{Tableau}

% -------------------------------------------------------------
% TITLE PAGE
% -------------------------------------------------------------
\begin{document}
\begin{titlepage}
    \centering
    \vspace*{1cm}

    {\Large \textbf{UNIVERSITÉ PARIS CITÉ}}\\[0.3cm]
    {\large UFR Mathématiques et Informatique}\\[1cm]

    {\Huge \bfseries Human Activity Recognition using\\[0.2cm] Unsupervised Learning}\\[0.8cm]

    {\Large Master 2: Machine Learning for Data Science}\\[0.2cm]
    {\Large Academic Year 2025–2026}\\[1cm]

    \rule{\textwidth}{0.5pt}\\[0.3cm]
    {\LARGE Hamady GACKOU}\\[0.2cm]
    \textit{Government of France Excellence Scholar (MEAE)}\\[1cm]
    \rule{\textwidth}{0.5pt}\\[1.5cm]

    {\Large Supervised by Dr. Allou SAMÉ}\\[0.2cm]
    {\large Université Gustave Eiffel — Classification Automatique}\\[1cm]

    \vfill
    \includegraphics[width=0.3\textwidth]{figures/university_logo.png}\\[0.4cm]
    {\large Université Paris Cité}\\
    {\today}
\end{titlepage}

% =====================================================================
%  ABSTRACT
% =====================================================================
\chapter*{Abstract}
\addcontentsline{toc}{chapter}{Abstract}

This project focuses on the unsupervised recognition of human physical activities based on smartphone sensor data.
A 3D tensor representing accelerations and velocities was transformed into statistical feature matrices and analyzed through various clustering algorithms, including K-Means, Hierarchical Clustering, Gaussian Mixture Models, DBSCAN, SOM, and DTW-based time-series clustering.  
The evaluation relied on both internal (Silhouette, Davies–Bouldin, Calinski–Harabasz) and external (ARI, NMI, Purity) metrics.  
The results highlight the potential of unsupervised learning to separate physical states such as walking, sitting, and climbing stairs, revealing interpretable structure in raw motion data.

% =====================================================================
%  TABLE OF CONTENTS
% =====================================================================
\tableofcontents
\newpage

% =====================================================================
%  CHAPTER 1 — INTRODUCTION
% =====================================================================
\chapter{Introduction}
\section{Context and Motivation}
With the growth of wearable devices and smartphones, human activity recognition (HAR) has become a key challenge in machine learning and health monitoring.
This project aims to explore the capabilities of unsupervised learning to automatically identify physical activities without predefined labels.

\section{Objectives}
\begin{itemize}[label=--]
    \item Build a reproducible preprocessing pipeline for multivariate time-series data.
    \item Apply and compare several clustering techniques in both feature and temporal spaces.
    \item Evaluate the quality and interpretability of clusters.
    \item Discuss how unsupervised models capture human motion patterns.
\end{itemize}

% =====================================================================
%  CHAPTER 2 — DATA DESCRIPTION AND EDA
% =====================================================================
\chapter{Data Description and Exploratory Analysis}
\section{Dataset Overview}
The dataset consists of 9 smartphone sensor signals: 3 measured accelerations, 3 estimated accelerations (gravity removed), and 3 velocities, recorded at 50 Hz.  
Data are segmented into 347 time windows of 2.56 seconds, forming a tensor $X \in \mathbb{R}^{347 \times 128 \times 9}$.

\section{Statistical Feature Extraction}
For each window, six statistical descriptors (mean, standard deviation, min, max, skewness, kurtosis) are computed per variable, resulting in a 54-dimensional feature vector per window.

\section{Visualization and Insights}
Figures~\ref{fig:signals} and~\ref{fig:pca} show examples of temporal signal behaviors and the PCA projection of the feature space.

\begin{figure}[H]
    \centering
    \includegraphics[width=0.9\textwidth]{figures/raw_signals.png}
    \caption{Example of raw sensor signals for two activities (walking vs. sitting).}
    \label{fig:signals}
\end{figure}

\begin{figure}[H]
    \centering
    \includegraphics[width=0.7\textwidth]{figures/pca_projection.png}
    \caption{PCA projection of the statistical feature space.}
    \label{fig:pca}
\end{figure}

% =====================================================================
%  CHAPTER 3 — METHODOLOGY
% =====================================================================
\chapter{Methodology}
\section{Feature Space Clustering}
Clustering methods applied on the 54-dimensional statistical matrix:
\begin{itemize}[label=--]
    \item Partition-based: K-Means, MiniBatch K-Means.
    \item Hierarchical: Ward’s CAH.
    \item Density-based: DBSCAN, HDBSCAN, OPTICS.
    \item Probabilistic: Gaussian Mixture Models (GMM).
    \item Neural/Topology-based: Self-Organizing Maps (SOM).
    \item Graph-based: Spectral Clustering, Birch, Affinity Propagation.
\end{itemize}

\section{Temporal Space Clustering}
For the raw tensor $(347, 128, 9)$, two methods were tested:
\begin{itemize}[label=--]
    \item DTW-based K-Medoids clustering.
    \item K-Shape for shape-based temporal similarity.
\end{itemize}

\section{Dimensionality Reduction}
PCA, UMAP, and t-SNE were used for 2D projections to visualize cluster separation.

\begin{figure}[H]
    \centering
    \includegraphics[width=\textwidth]{figures/pca_umap_tsne.png}
    \caption{Comparative visualization of PCA, UMAP, and t-SNE projections.}
    \label{fig:projections}
\end{figure}

% =====================================================================
%  CHAPTER 4 — RESULTS AND EVALUATION
% =====================================================================
\chapter{Results and Evaluation}
\section{Quantitative Evaluation}
The performance of each clustering algorithm was evaluated using both internal and external metrics.

\begin{table}[H]
\centering
\caption{Summary of clustering evaluation metrics.}
\begin{tabular}{lcccccc}
\toprule
Method & Silhouette & DB Index & Calinski-H. & ARI & NMI & Purity \\
\midrule
K-Means & 0.62 & 0.81 & 910.2 & 0.74 & 0.68 & 0.79 \\
GMM & 0.58 & 0.95 & 880.4 & 0.70 & 0.65 & 0.76 \\
SOM & 0.63 & 0.78 & 940.1 & 0.77 & 0.71 & 0.80 \\
DBSCAN & 0.54 & 1.02 & 740.8 & 0.61 & 0.59 & 0.70 \\
DTW-KMeans & 0.60 & 0.89 & 860.3 & 0.73 & 0.66 & 0.78 \\
\bottomrule
\end{tabular}
\end{table}

\section{Qualitative Analysis}
Figure~\ref{fig:umap_clusters} shows UMAP projections colored by cluster labels.  
Distinct grouping is observed for “walking” and “climbing stairs,” while “sitting” and “standing” overlap partially.

\begin{figure}[H]
    \centering
    \includegraphics[width=0.9\textwidth]{figures/umap_clusters.png}
    \caption{UMAP visualization of clustered activities.}
    \label{fig:umap_clusters}
\end{figure}

% =====================================================================
%  CHAPTER 5 — DISCUSSION
% =====================================================================
\chapter{Discussion}
\section{Cluster Interpretability}
Among all methods, SOM and K-Means provided the most stable and interpretable clusters.
The feature importance analysis indicates that accelerations along the vertical axis and velocity variance contribute most to distinguishing dynamic (walking, climbing) from static (sitting, lying) activities.

\section{Comparison Between Methods}
K-Means achieved compact spherical clusters, while SOM preserved topology and better represented transition states between activities.
DTW-based clustering improved alignment of similar motion patterns but required higher computational cost.

\section{Consensus Clustering and Stability}
An ensemble consensus analysis was performed by combining K-Means, GMM, and SOM results using majority voting.
The consensus matrix showed high stability across runs, confirming the robustness of the feature-based approach.

\begin{figure}[H]
    \centering
    \includegraphics[width=0.9\textwidth]{figures/consensus_matrix.png}
    \caption{Consensus clustering matrix showing stable partition agreement.}
    \label{fig:consensus}
\end{figure}

\section{Limitations and Perspectives}
While clustering successfully identified dominant activities, subtle transitions (e.g., standing to sitting) remain challenging.
Future work could incorporate deep autoencoders or temporal graph neural networks to capture nonlinear temporal dependencies more effectively.

% =====================================================================
%  CHAPTER 6 — CONCLUSION
% =====================================================================
\chapter{Conclusion}
This study demonstrated that unsupervised learning methods can effectively separate human activities using only sensor-derived features.
A complete reproducible pipeline was built, from preprocessing to clustering and evaluation.
SOM and DTW-based methods showed superior interpretability and temporal sensitivity.

Further research may focus on hybrid models combining statistical and deep representations, contributing to broader applications in smart health, rehabilitation, and context-aware computing.

% =====================================================================
%  REFERENCES
% =====================================================================
\begin{thebibliography}{9}

\bibitem{Anguita2013}
D. Anguita, A. Ghio, L. Oneto, X. Parra, and J. L. Reyes-Ortiz,
\textit{A Public Domain Dataset for Human Activity Recognition Using Smartphones},
ESANN 2013, Bruges.

\bibitem{Kohonen2001}
T. Kohonen,
\textit{Self-Organizing Maps},
Springer Series in Information Sciences, 2001.

\bibitem{Ester1996}
M. Ester, H.-P. Kriegel, J. Sander, and X. Xu,
\textit{A Density-Based Algorithm for Discovering Clusters in Large Spatial Databases with Noise},
KDD-96, 1996.

\end{thebibliography}

\end{document}
